% \vspace{-2mm}
\section{\method}
% \vspace{-2mm}
% We introduce \method (\textbf{S}ynthetic yet \textbf{U}nbiased \textbf{G}eneration with \textbf{A}dversarially \textbf{R}efined \textbf{CREPE}), a new benchmark for faithful evaluation of VLMs' compositionality.
% \method presents two key contributions over existing benchmarks: (1) it drastically reduces the two identified dataset biases, and (2) it covers a broad range of fine-grained types of hard negatives. We present a summary comparison on compositionality benchmarks in Appendix~\ref{app:bms}.
% In particular, \method is built upon the image-text pairs of COCO~\cite{lin2014microsoft} and its hard negative generation pipeline consists of two main stages (Sec.~\ref{sec:data_generation}). First, we generate hard negative texts by prompting ChatGPT~\cite{chatgpt} with few-shot demonstrations to modify a given positive text into corresponding hard negatives, since it is capable of generating sensical and fluent texts.
% Second, we employ an adversarial refinement mechanism to maximally disentangle nonsensical and non-fluent biases from \method, avoiding it from being hacked via the biases.
% We apply this data generation pipeline to a broad range of fine-grained types of hard negatives (Sec.~\ref{sec:hn_types}).

We introduce \method, a new benchmark for faithful evaluation of vision-language models' compositionality based on the image-text pairs of COCO~\cite{lin2014microsoft}.
\method presents two key contributions over existing benchmarks: (1) it drastically reduces the two identified dataset biases (Sec.~\ref{sec:data_generation}), and (2) it covers a broad range of fine-grained types of hard negatives (Sec.~\ref{sec:hn_types}). We present a summary comparison on compositionality benchmarks in Appendix~\ref{app:bms}.
% In particular, \method is built upon the image-text pairs of COCO~\cite{lin2014microsoft} and its hard negative generation pipeline consists of two main stages (Sec.~\ref{sec:data_generation}). First, we generate hard negative texts by prompting ChatGPT~\cite{chatgpt} with few-shot demonstrations to modify a given positive text into corresponding hard negatives, since it is capable of generating sensical and fluent texts.
% Second, we employ an adversarial refinement mechanism to maximally disentangle nonsensical and non-fluent biases from \method, avoiding it from being hacked via the biases.
% We apply this data generation pipeline to a broad range of fine-grained types of hard negatives (Sec.~\ref{sec:hn_types}).


% nonfluent and nonsensical candidates by leveraging both ChatGPT's general knowledge and adversarial filtering techniques~\cite{}.
% Finally, we crowdsource human annotators to look at both the images and the generated hard negatives to ensure the data validity.



\subsection{\method generation workflow alleviates dataset biases}
\label{sec:data_generation}

% Existing benchmarks are hackable due to the fact that for a significantly \emph{unbalanced portion} of the dataset, the hard negative texts are \emph{less} sensical and fluent than the positive texts.
% In response, \method generates more sensical and fluent hard negatives to close the distributional gaps to the positive texts. Then, it ensures a balanced distribution on the commonsense and fluency score gaps to make the final evaluation dataset informationless and thus robust w.r.t. to the identified biases.
The generation procedure of \method consists of three main stages, centered around creating sensical and fluent hard negatives that close the distributional gaps to the positive texts, and ensuring a balanced distribution on the score gaps to make the final dataset robust to the identified biases.

\textbf{Stage 1: Generate sensical and fluent hard negatives with a large language model.}
Observing the capability of modern large language models in generating fluent and plausible texts, we leverage ChatGPT~\cite{chatgpt} to generate hard negative texts where we explicitly instruct it to avoid commonsense (logical) and fluency (grammatical) errors. To guide ChatGPT in re-writing a given positive text into its hard negative counterparts, we provide few-shot demonstrations written by the authors and leverage its in-context learning ability to generalize to unseen texts.
Figure~\ref{fig:example_prompt} shows an example demonstration used and an actual hard negative generated.
We detail all the prompt templates in Appendix~\ref{app:prompt}.
Table~\ref{tab:sugar_crepe_examples} shows the comparisons between hard negatives generated from ChatGPT in \method and that from existing benchmarks.

\textbf{Stage 2: Filter false negatives with human validation.}
A generated text is considered a valid hard negative only if it incorrectly describes the corresponding image. For example, given an image with a positive caption ``a man and a child sitting on a sofa'', a compositional change that replaces ``child'' with ``girl'' may still result in a correct caption.
To ensure the validity of the hard negatives in \method, we filter out \textit{false} negatives by manually examining the generated hard negatives and their corresponding images.

\begin{figure}
    \centering
    \begin{minipage}{0.52\textwidth}
        \centering
        \includegraphics[width=\linewidth]{imgs/sugar-crepe-prompt.pdf}
        \caption{Example prompt (black) and actual hard negative (green) generated from  ChatGPT.}
        \label{fig:example_prompt}
    \end{minipage}\hfill
    \begin{minipage}{0.46\textwidth}
\vspace{-3mm}
\begin{algorithm}[H]  
% \scalebox{0.65}{
\small
  \caption{Adversarial Refinement}
  \label{alg:ar}  
  \begin{algorithmic}[1]  
   \Require 
      Text-only model $M_1$ and $M_2$; Number of grids $K$;
      A set of candidates $\mathcal{D}=\left\{I_i, T^{\mathrm{p}}_i, T^{\mathrm{n}}_i\right\}_{i\in [N]}$, where $I_i$, $T^{\mathrm{p}}_i$, and $T^{\mathrm{n}}_i$ are $i$-th image, positive caption, and negative caption.
   \Ensure A subset $\bar{\mathcal{D}}\subset\mathcal{D}$
   \State Calculate the model score gap for each candidate $g^{(1)}_i=M_1(T^\mathrm{p}_i)-M_1(T^\mathrm{n}_i)$ and $g^{(2)}_i=M_2(T^\mathrm{p}_i)-M_2(T^\mathrm{n}_i)$
   \State Split the 2D space $[-1, 1]\times[-1, 1]$ to $K\times K$ equal-size grids.
   \State Place each candidate to a grid based on the score gaps $g^{(1)}_i$ and $g^{(2)}_i$.
   \State Initialize $\bar{\mathcal{D}}=\{\}$
    \For{each pair of grid $(G_j, G^{*}_j)$ symmetric about the original point $(0, 0)$}
        \If{$|G_j| > |G^*_j|$}
        \State Sample $|G^*_j|$ candidates from $G_j$ and put them to $\bar{\mathcal{D}}$. 
        \State Put candidates in $G^*_j$ to $\bar{\mathcal{D}}$. 
        \Else
        \State Sample $|G_j|$ candidates from $G^*_j$ and put them to $\bar{\mathcal{D}}$. 
        \State Put candidates in $G_j$ to $\bar{\mathcal{D}}$. 
        \EndIf
    \EndFor 
  \end{algorithmic}  
% }
\end{algorithm}  
    \end{minipage}
\vspace{-2ex}
\end{figure}



\textbf{Stage 3: De-bias dataset with adversarial refinement.}
While ChatGPT yields more sensical and fluent text, there is no guarantee that the bias between positive and negative texts is negligible.
Following dataset de-biasing work~\cite{zellers2018swag,sakaguchi2021winogrande,le2020adversarial}, we develop an adversarial refinement mechanism that maximally reduces the undesirably exploitable artifacts in \method.
Specifically, our goal is to ensure that performance improvements on \method cannot be achieved by exploiting the identified nonsensical and non-fluent biases. To accomplish this, we characterize the biases again with the commonsense and grammar models~\cite{liu2023vera,morris2020textattack}, and subsample the dataset to ensure symmetric score gap distributions on both the positive and negative sides, as shown in Figure~\ref{fig:vera_gap_ours}. 
We note the symmetry around zero implies that the commonsense and grammar scores can no longer be used to infer the ground truth positive texts.
% this is effectively minimizing the mutual information available to infer the ground truth positive texts solely from the score gaps.
We provide the adversarial refinement algorithm in Algorithm~\ref{alg:ar}.


% \begin{wrapfigure}[22]{R}{0.48\textwidth}
% \vspace{-2\intextsep}
%  \begin{minipage}{0.5\textwidth}
% \begin{algorithm}[H]  
% % \scalebox{0.65}{
% \small
%   \caption{Adversarial Refinement}
%   \label{alg:ar}  
%   \begin{algorithmic}[1]  
%    \Require 
%       Text-only model $M_1$ and $M_2$; Number of grids $K$;
%       A set of candidates $\mathcal{D}=\left\{I_i, p_i, n_i\right\}_{i\in [N]}$, where $I_i$, $p_i$, and $n_i$ are $i$-th image, positive caption, and negative caption.
%    \Ensure A subset $\bar{\mathcal{D}}\subset\mathcal{D}$
%    \State Calculate the model score gap for each candidate $g^{(1)}_i=M_1(p_i)-M_1(n_i)$ and $g^{(2)}_i=M_2(p_i)-M_2(n_i)$
%    \State Split the 2D space $[-1, 1]\times[-1, 1]$ to $K\times K$ equal-size grids.
%    \State Place each candidate to one grid based on its score gap $g^{(1)}_i$ and $g^{(2)}_i$.
%    \State Initialize $\bar{\mathcal{D}}=\{\}$
%     \For{each pair of grid $(G_j, G^{*}_j)$ symmetric about the original point $(0, 0)$}
%         \If{$|G_j| > |G^*_j|$}
%         \State Sample $|G^*_j|$ candidates from $G_j$ and put them to $\bar{\mathcal{D}}$. 
%         \State Put candidates in $G^*_j$ to $\bar{\mathcal{D}}$. 
%         \Else
%         \State Sample $|G_j|$ candidates from $G^*_j$ and put them to $\bar{\mathcal{D}}$. 
%         \State Put candidates in $G_j$ to $\bar{\mathcal{D}}$. 
%         % \EndIf
%     % \EndFor 
%   \end{algorithmic}  
% % }
% \end{algorithm}  
%  \end{minipage}
%   \end{wrapfigure}

\subsection{\method covers a broad range of hard negative types}
\label{sec:hn_types}

To test different aspects of vision-language models' compositional understanding, we follow CREPE~\cite{ma2022crepe} to consider various \emph{forms} of hard negatives, and follow VL-CheckList~\cite{zhao2022vl} and ARO~\cite{yuksekgonul2023when} to consider different fine-grained \emph{categories} of the atomic concepts.
In total, \method covers $7$ fine-grained types of hard negatives, as shown in Table~\ref{tab:dataset_summary}.
We introduce the dataset taxonomy below, starting from the \emph{form} of the hard negatives to its different \emph{finer-grained} variants.



\textbf{The \replace\ form.}
Given a positive text describing a scene, we generate a \replace\ hard negative by replacing an atomic concept in the original text with a new concept that makes the text mismatch with the original scene.
Based on the type of the atomic concept---object, attribute, or relation---we further categorize \replace\ hard negatives into \replaceobj, \replaceatt, and \replacerel.

\textbf{The \swap\ form.} Different from \replace, \swap\ does not introduce new concepts in the hard negatives, but a \swap\ hard negative is generated by swapping two atomic concepts of the same category in the positive text. We further categorize \swap\ into \swapobj\ and \swapatt, and omit swapping two relationships since it generally results in nonsensical texts.

\textbf{The \add\ form.} Similar to the \replace\ form, but instead of replacing an atomic concept with a new one, we generate an \add\ hard negative by adding a new atomic concept to the positive text that makes it mismatch with the original scene.
% Note that \add\ hard negatives are designed to test a model's ability follows the spirit are intended to probe a VLM's sensitivity to 
We only further categorize \add\ into \addobj~(adding object concept) and \addatt~(adding attribute concept), as adding new relationship concepts to the positive texts often make them highly implausible.

% \subsection{\method employs human validation for hard negatives}
% To ensure the generated hard negatives are valid, \ie, the compositional change indeed make the. \method utilizes manual examination to filter out invalid hard negatives.

\textbf{Dataset overview.} The final evaluation set of \method consists of $7512$ examples, where the numbers for each fine-grained type are listed in Table~\ref{tab:dataset_summary}.
Each example is an image-to-text retrieval task composed of an image, a positive text, and a hard negative. On \method, random chance performance has an average accuracy of $50\%$.
We note that ARO and CREPE additionally consider \shuffle\ (randomly shuffling words in a sentence) and \negate\ (adding negation keywords ``no/not'' to a sentence) hard negatives. We however omit them in \method as \shuffle\ is very unlikely to be plausible and fluent, and \negate\ introduces irreducible keyword artifacts~\cite{ma2022crepe}.~\footnote{One can easily infer hard negatives from whether the text contains negation keywords ``no/not''.}


\begin{table*}[!t]
  \centering
  \small
  \caption{We report the number of hard negative captions of all types in \method.}
  \scalebox{0.7}{
    \begin{tabular}{lccccccc} 
    \toprule
     & \multicolumn{3}{c}{\bf \replace} & \multicolumn{2}{c}{\bf \swap} &  \multicolumn{2}{c}{\bf \add}\\ 
     \cmidrule(lr){2-4}\cmidrule(lr){5-6}\cmidrule(lr){7-8}
     
     & \textbf{Object} & \textbf{Attribute} & \textbf{Relation} &\textbf{Object} & \textbf{Attribute} & \textbf{Object} & \textbf{Attribute} \\ 
    \midrule

\textbf{\# negative captions} & 1652 & 788 & 1406  &  246 & 666  & 2062 & 692 \\

\bottomrule 
\end{tabular}
}
  
  \label{tab:dataset_summary}
\end{table*}